\section{动量与 Adam 在平均什么}
\label{sec:14-Deep-Learning/03-Optimization-and-Training/02}

SGD 的每一步只看当前 mini-batch 的梯度：

\[
\theta_t=\theta_{t-1}-\eta_t g_t,
\qquad
g_t\approx\nabla J(\theta_{t-1}).
\]

\(g_t\) 同时被两股力量扭曲：目标函数的局部曲率，以及 mini-batch 采样引入的随机噪声。在一个狭长的谷底中，梯度在陡峭方向上来回换号——这步让你往左，下步又把你拉回来；在平缓方向却持续同号——每步都推着参数沿谷底慢慢前进。mini-batch 噪声叠加在这两个方向上，让每一步都额外抖动一次。

把这种情况推到极端，能看得更清楚。假设一个二维损失函数 \(J(w_1,w_2)=0.1w_1^2+10w_2^2\)，等高线是极度拉长的椭圆：\(w_2\) 方向的曲率是 \(w_1\) 方向的 100 倍。梯度为 \(\nabla J=(0.2w_1,\;20w_2)\)。在点 \((5,0.1)\) 处，梯度是 \((1,\;2)\)——\(w_2\) 方向的梯度反而更大，即使 \(w_2\) 已经离原点很近。SGD 会沿 \((-1,-2)\) 方向迈步，在 \(w_2\) 方向上跨过原点、翻到对侧——下一个梯度又指向回来。如果学习率稍大，便是一场永不收敛的振荡。而在 \(w_1\) 方向，梯度只有 1，参数从 5 走到 0 需要很多步。

一个被陡峭方向折腾得很惨的 SGD，和一个在平缓方向走不动的 SGD，都渴望同样的东西：从历史梯度中过滤掉高频换号的分量，保留持续同号的趋势。动量和 Adam 做的正是这件事，但它们平均的对象不同。

\subsection{动量积累的是方向的一阶时间平均}

动量的核心操作：维护一个移动平均 \(v_t\)，每一步用当前梯度更新它，再用它（不是裸梯度）来更新参数。一种常见写法：

\[
v_t=\beta v_{t-1}+(1-\beta)g_t,
\qquad
\theta_t=\theta_{t-1}-\eta v_t,
\]

其中 \(0\le\beta<1\)。展开递推（设 \(v_0=0\)）：

\[
v_t=(1-\beta)\sum_{k=1}^{t}\beta^{t-k}g_k.
\]

这不是等权平均。最近一步梯度权重为 \(1-\beta\)，前一步为 \((1-\beta)\beta\)，再往前为 \((1-\beta)\beta^2\)——按几何速度衰减。有效记忆长度的量级约为 \(1/(1-\beta)\)：\(\beta=0.9\) 时约 10 步，\(\beta=0.99\) 时约 100 步。\(\beta\) 越大，平均越平滑，但对最新梯度的响应越迟钝——和平滑信号的永恒取舍。

回到狭长椭圆。\(w_2\) 方向上，梯度步与步之间反复换号：\(g_1=+2.0\), \(g_2=-1.6\), \(g_3=+1.3\), \(g_4=-1.0\)。取 \(\beta=0.9\)，按公式递推几步：

\[
\begin{aligned}
v_1&=0.1\times 2.0=0.200,\\
v_2&=0.9\times 0.200+0.1\times(-1.6)=0.180-0.160=0.020,\\
v_3&=0.9\times 0.020+0.1\times 1.3=0.018+0.130=0.148,\\
v_4&=0.9\times 0.148+0.1\times(-1.0)=0.133-0.100=0.033.
\end{aligned}
\]

\(v_t\) 在正负交替中互相抵消，绝对值始终在零点附近小幅度振荡——陡峭方向的移动被抑制了。

\(w_1\) 方向上，每步梯度同号：\(g_1=+1.0\), \(g_2=+0.8\), \(g_3=+0.6\), \(g_4=+0.5\)。同样的 \(\beta=0.9\)：
\[
\begin{aligned}
v_1&=0.100,\\
v_2&=0.090+0.080=0.170,\\
v_3&=0.153+0.060=0.213,\\
v_4&=0.192+0.050=0.242.
\end{aligned}
\]
\(v_t\) 持续在同方向积累。四步之后，动量在 \(w_1\) 方向的累积值（0.242）已远超 \(w_2\) 方向的残余（0.033），尽管每一步 \(w_2\) 方向的裸梯度绝对值都更大。参数沿谷底滑行的速度被动量放大，同时不会因陡峭方向的振荡被推出谷底。

有些资料使用 \(v_t=\beta v_{t-1}+g_t\)，把 \(1-\beta\) 因子吸收进学习率。两套参数的 \(\eta\) 在数值上不相等——比较实验或复现结果时必须先统一记号，不能只看 \(\eta\) 数值。

训练初期 \(v_t\) 从零开始累积，权重和是 \(1-\beta^t<1\)，会系统性地偏小。如果需要把 \(v_t\) 解释为梯度均值，可以用偏差修正 \(\widehat v_t=v_t/(1-\beta^t)\)。普通 momentum 更新不一定做显式修正——学习率和 warmup 也会吸收初期的尺度偏差——但知道这个偏小效应存在，有助于理解为什么训练前几步的有效步长总是偏小。

\subsection{Adam 同时估计一阶尺度与二阶尺度}

动量解决了一个方向问题：在振荡维度上自动抵消，在持续维度上自动积累。但它留下了一个尺度问题：不同参数的梯度量级可以差两三个数量级。

考虑一个同时包含嵌入层和顶层分类器的网络。嵌入层只有当前 batch 中出现的少数 token 会收到非零梯度，这些梯度往往数值较大（因为稀疏更新缺少平均稀释）；顶层分类器的每个参数每条样本都收到梯度，数值稳定但偏小。单一全局学习率无法同时伺候：对嵌入层恰好合适的学习率，对顶层分类器太小；对顶层合适的，嵌入层被炸飞。

Adam 在动量之上叠加了第二个平均——梯度平方的指数移动平均，用来逐坐标估计梯度的典型大小。

维护两个移动平均：
\[
\begin{aligned}
m_t&=\beta_1m_{t-1}+(1-\beta_1)g_t,\\
v_t&=\beta_2v_{t-1}+(1-\beta_2)g_t\odot g_t.
\end{aligned}
\]
\(m_t\) 是动量——梯度方向的时间平均，与上一节做的事相同。\(v_t\) 是梯度平方的时间平均——每个坐标上梯度有多``活跃''。\(\beta_1\) 典型取 0.9，\(\beta_2\) 取 0.999：二阶矩的衰减比一阶慢得多，因为梯度平方的波动本身就比梯度方向慢——一个参数的梯度的绝对值通常不会在几步之内突变一个数量级。

偏差修正后：
\[
\widehat m_t=\frac{m_t}{1-\beta_1^t},
\qquad
\widehat v_t=\frac{v_t}{1-\beta_2^t},
\]
更新为逐坐标的除法与开方：
\[
\theta_t
=\theta_{t-1}
-\eta\frac{\widehat m_t}{\sqrt{\widehat v_t}+\varepsilon}.
\]

分子 \(\widehat m_t\) 是持续方向。分母 \(\sqrt{\widehat v_t}\) 是近期梯度的 RMS——一个坐标的梯度长期较大，说明这个方向敏感，有效步长应当收敛；梯度稀疏或微弱，说明这个方向不活跃，有效步长应当相对放大。注意 \(\sqrt{\widehat v_t}\) 依赖的是梯度平方的历史，也就是依赖训练轨迹——它不是目标函数在该坐标方向的固有曲率，而是``优化器最近在这个方向上经历了多大的梯度''。换一个初始化点、换一个 batch 顺序，\(v_t\) 可能完全不同。

\(\varepsilon\) 通常在 \(10^{-8}\) 量级，正式作用是除零保护。当 \(\sqrt{\widehat v_t}\) 接近 \(\varepsilon\) 时（训练初期的很多坐标），\(\varepsilon\) 会显著改变有效步长——这时分母主要是 \(\varepsilon\)，分子很小，有效步长反而最大。在 mixed precision 训练中，float16 的最小正数又改变了``接近零''的边界含义。因此 \(\varepsilon\) 是一个有实际影响的算法参数。

\subsection{自适应缩放不等于 Hessian 的逆}

读到除法这一步，你可能联想到 Newton 法：Newton 法用 Hessian 的逆 \(H^{-1}\) 对梯度做线性变换，使更新方向沿损失景观的曲率等值面径直指向极小值。Adam 的逐坐标除法也是缩放了梯度——但它和 Newton 法之间的差距不能忽略。

关键差距有三层。第一，\(v_t\) 来自梯度平方的 EMA，不是目标函数的二阶导数——梯度很大不一定意味着曲率很大，也可能只是当前参数恰好在一个陡坡上。第二，逐坐标除法只提供对角缩放，完全忽略参数之间的相关方向。回到那个旋转后的狭长椭圆：如果曲率的主方向不是沿坐标轴方向，而是沿 \((1,1)\) 方向，对角缩放对此毫无知觉，Adam 仍会在谷底走出之字形。第三，\(v_t\) 依赖训练轨迹——换一个初始化或 batch 顺序，\(v_t\) 就不同，有效步长随之改变。Newton 法的 Hessian 是目标函数的固有属性，与优化路径无关。

既然如此，Adam 的优势在哪里？在那些坐标方向的尺度差异天然就很大的问题上——稀疏特征、嵌入参数、不同层之间的激活量级——逐坐标适应提供的初期加速是真实而显著的。在梯度稀疏到几乎为零的坐标上，动量的一阶平均也可能几乎为零，但二阶平均 \(v_t\) 同样很低，有效步长因此被放大，使得这些参数不会完全冻结。SGD + momentum 在同等学习率下，稀疏坐标可能要上百步才能收到一次非零梯度，期间参数纹丝不动。

这解释了 Adam 在 NLP 和 Transformer 训练中的普及——不是因为 Adam 在所有目标上优于 SGD 的定理，而是因为它以逐坐标适应的方式处理了参数间天然的尺度差异，降低了调参的初期门槛。

\subsection{权重衰减与 \(L_2\) 正则在自适应方法中分道扬镳}

在 SGD 中，\(L_2\) 正则和权重衰减是同一操作的两种命名。若目标函数加入 \(\frac{\lambda}{2}\norm{\theta}^2\)，梯度变为 \(g_t+\lambda\theta_{t-1}\)。SGD 更新：
\[
\theta_t=\theta_{t-1}-\eta(g_t+\lambda\theta_{t-1})
=(1-\eta\lambda)\theta_{t-1}-\eta g_t.
\]
每个坐标被同等比例地收缩——\(L_2\) 正则等价于每步将参数向原点缩放 \(1-\eta\lambda\)。

但在 Adam 中，正则化梯度同样被逐坐标除以 \(\sqrt{\widehat v_t}+\varepsilon\)：
\[
\theta_t=\theta_{t-1}-\eta\frac{\widehat m_t+\lambda\theta_{t-1}}{\sqrt{\widehat v_t}+\varepsilon}.
\]
梯度大的坐标：分母大，步长小——不仅梯度更新被缩小，正则化收缩也被缩小。梯度小的坐标：分母小，步长大——收缩反而强。同一个 \(\lambda\)，在不同坐标上的实际收缩力度从此不再相同。

AdamW 明确把衰减和自适应梯度拆成两项：
\[
\theta_t
=(1-\eta\lambda)\theta_{t-1}
-\eta\frac{\widehat m_t}{\sqrt{\widehat v_t}+\varepsilon}.
\]
前半部分：对所有参数同等比例的权重收缩，与梯度无关。后半部分：自适应梯度更新。两者优化的是不同的离散动力学——Adam 的 \(L_2\) 正则是``目标和优化器耦合的收缩''，AdamW 的权重衰减是``优化器外的独立收缩''。

选择哪一种，需要说明你希望表达的是目标函数中的参数惩罚（Adam + \(L_2\)），还是优化过程中与梯度信号无关的直接收缩（AdamW）。两者在不同任务上的相对优劣，取决于哪种收缩方式更好地配合了自适应缩放——目前没有统一的理论结论，实践中 AdamW 被更广泛地使用，部分原因是它的收缩不因坐标活跃度而变，行为更可预测。

\subsection{优化器不是训练的万能开关}

Adam 能快速把训练 loss 压下去，不说明验证集上更好。SGD 在后期走得慢而稳定，也不说明它一定拥有更好的隐式正则化。比较优化器时必须控制学习率调度、batch size、总训练步数、权重衰减的定义方式和初始化——然后同时观察训练目标和独立验证结果。

训练失败时的诊断优先级：先确认几十个样本能否被过拟合（排错 shape/loss/gradient）；再检查梯度中是否有 inf/NaN（排错数值稳定）；再计算参数更新量相对参数本身尺度的比值（排错步长）；再查各层激活和梯度是否跨层爆炸或消失（排错初始化/归一化）。这些检查没通过之前，直接把 SGD 换成 Adam 可能暂时藏住尺度问题，却不会修复错误标签、错误 reduction、数据泄漏或不可识别的目标定义。

动量和 Adam 的共同价值：把单步噪声梯度变成带记忆、带尺度感知的更新，让优化器在陡峭方向自己刹车，在平缓方向自己加速。理解它们在平均什么、怎样改变每个坐标的有效步长，才能把优化器当作可分析、可诊断的数值工具——而不是训练脚本里一个不曾细读的默认参数。

本节解释更新量的组成；\hyperref[sec:09-Convex-Optimization/08-Nonconvex-and-Stochastic-Optimization/05]{``动量与自适应缩放改变的是动力学''}进一步把 momentum 写成曲率方向上的二阶递推，区分稳定、振荡与发散三种动力学模式，并说明 Adam 为什么不能从随机梯度的无偏性直接继承收敛保证。
